\documentclass[11pt]{article}

% -------------------------
% Packages
% -------------------------
\usepackage{amsmath, amssymb, amsthm, mathtools}
\usepackage{geometry}
\usepackage{tikz}
\usepackage{hyperref}
\usepackage{enumitem}

\geometry{margin=1in}

% -------------------------
% Theorems
% -------------------------
\newtheorem{theorem}{Theorem}[section]
\newtheorem{definition}[theorem]{Definition}
\newtheorem{proposition}[theorem]{Proposition}
\newtheorem{lemma}[theorem]{Lemma}
\newtheorem{remark}[theorem]{Remark}
\newtheorem{axiom}[theorem]{Axiom}

% -------------------------
% Title
% -------------------------
\title{\textbf{The Coh Category:\\
Verifier-Governed Admissibility and Constraint-Enriched Structure}}
\author{NoeticanLabs Research Division}
\date{April 23, 2026}

\begin{document}

\maketitle

% =====================================================
\begin{abstract}
We introduce the Coh category, a verifier-indexed, constraint-enriched categorical structure in which morphisms are generated by admissibility rather than assumed a priori. 

We prove that verifier-certified transitions form a category under telescoping inequality closure, and that projection-induced information loss yields an oplax composition law.

We further demonstrate that bicategorical structure arises naturally from slack in composition, where 2-cells explicitly represent the optimization gap between syntactic operational ledgers and semantic physical bounds. 

The framework provides a rigorous foundation for governed computation, explicitly separating observable dynamics from hidden state evolution through a locally posetal 2-category.
\end{abstract}

% =====================================================
\section{Introduction}

Traditional category theory assumes morphisms exist independently of constraints. In real systems, transitions must satisfy verification and resource constraints. 

\[
\boxed{
\text{Morphisms exist if and only if they are verifier-admissible}
}
\]

Unlike enriched categories, where morphisms are decorated with costs or distances after construction, the Coh framework generates the morphisms from the admissibility verification itself. This isolates physical constraints into a verifiable ledger.

% =====================================================
\section{Coh Systems}

\begin{definition}[Coh System]
A Coh System is a tuple 
\[
\mathcal{S} = (X, \mathcal{K}, V, \Pi, \mathrm{RV})
\]
where:
\begin{itemize}
\item $X$ is the physical or computational state space.
\item $\mathcal{K} \subseteq X$ is the closed, convex set of admissible states.
\item $V: X \to \mathbb{R}_{\ge 0} \cup \{+\infty\}$ is a proper, lower semicontinuous violation functional.
\item $\Pi$ is the macroscopic observation projection.
\item $\mathrm{RV}$ is the deterministic receipt verifier.
\end{itemize}
\end{definition}

\subsection{Hidden State and Fiber Geometry}

\begin{definition}[Hidden State]
Each observable macroscopic state $x$ corresponds to a fiber of hidden microscopic encodings $h_x \in \mathcal{H}$, such that the projection $\Pi(h_x) = x$.
\end{definition}

% =====================================================
\section{Admissibility and The Verifier}

\begin{definition}[Admissible Transition]
A transition morphism $r: x \to x'$ exists if and only if the receipt verifier accepts the proposed trace,
\[
\mathrm{RV}(x, r, x') = \mathrm{ACCEPT}
\]
and the strict thermodynamic bounds are satisfied:
\[
V(x') + s(r) \le V(x) + d(r)
\]
where $s(r)$ is the macroscopic spend (work done) and $d(r)$ is the allowed defect (unpaid slack).
\end{definition}

\subsection{Verifier as Structure}

\begin{proposition}
The verifier $\mathrm{RV}$ induces a strict predicate defining a partial subcategory of valid transitions over $X$.
\end{proposition}
\begin{proof}
Let $\mathcal{G}$ be the complete directed graph of all conceivable transitions on $X$. The verifier $\mathrm{RV}$ acts as a boolean filter $\mathcal{G} \to \{0,1\}$. The subgraph induced by the preimage $\mathrm{RV}^{-1}(1)$ defines the existence of edges. Therefore, the morphisms of the Coh category are exactly the elements of this filtered subgraph, making the verifier fundamentally structural.
\end{proof}

% =====================================================
\section{Projection and The Oplax Envelope}

To bridge continuous unobserved physics and discrete recorded ledgers, we analyze the information loss across the projection $\Pi$.

\begin{definition}[Fiber of a Record]
For a macroscopic transition record $R = \Pi(\Psi)$, the unobserved fiber is:
\[
\mathcal{F}_R = \{\Psi : \Pi(\Psi)=R\}
\]
\end{definition}

\begin{definition}[Envelope Cost]
The minimal safe macroscopic spend $s(R)$ bounding the true continuous work $W(\Psi)$ must satisfy the supremum over the hidden fiber:
\[
s(R) \ge \sup_{\Psi \in \mathcal{F}_R} W(\Psi)
\]
\end{definition}

\begin{proposition}[Oplax Bound]
The macroscopic spend functional is strictly subadditive (oplax) under composition:
\[
s(r_2 \circ r_1) \le s(r_1)+s(r_2)
\]
\end{proposition}
\begin{proof}
Let $R_1$ and $R_2$ be sequentially composable macroscopic records. The true continuous work functional $W$ is additive over time, such that for any continuous path $\Psi \in \mathcal{F}_{R_2 \circ R_1}$ spanning both intervals, $\Psi$ can be split into $\Psi_1 \in \mathcal{F}_{R_1}$ and $\Psi_2 \in \mathcal{F}_{R_2}$ where $W(\Psi) = W(\Psi_1) + W(\Psi_2)$. 

Taking the supremum over the composite fiber:
\begin{align*}
\sup_{\Psi \in \mathcal{F}_{R_2 \circ R_1}} W(\Psi) &= \sup_{\Psi_1, \Psi_2} [W(\Psi_1) + W(\Psi_2)] \\
&\le \sup_{\Psi_1 \in \mathcal{F}_{R_1}} W(\Psi_1) + \sup_{\Psi_2 \in \mathcal{F}_{R_2}} W(\Psi_2) \\
&\le s(r_1) + s(r_2)
\end{align*}
Thus, evaluating the envelope cost $s$ across the concatenated record yields $s(r_2 \circ r_1) \le s(r_1) + s(r_2)$. The loss of internal boundary conditions forces the inequality.
\end{proof}

% =====================================================
\section{The Coh Category}

\begin{definition}[Identity]
For any $x \in X$, the identity morphism $\mathrm{id}_x$ is defined by the zero-receipt, satisfying:
\[
s(\mathrm{id}_x)=0,\quad d(\mathrm{id}_x)=0
\]
\end{definition}

\subsection{Composition Semantics}

\subsubsection{Syntactic Composition}
Syntactic composition represents the raw concatenation of receipts, preserving strict associativity:
\[
s_{\mathrm{syn}} = s_1 + s_2, \quad d_{\mathrm{syn}} = d_1 + d_2
\]

\subsubsection{Semantic Composition}
Semantic composition represents the evaluated envelope cost of the combined interval. Due to projection-induced information loss, the semantic bound is typically tighter than the syntactic sum:
\[
s_{\mathrm{sem}} \le s_{\mathrm{syn}}
\]

\begin{theorem}
$\mathbf{Coh}(\mathcal{S})$ forms a category under verifier-certified composition.
\end{theorem}
\begin{proof}
Objects are the states $X$. Morphisms are the verifier-accepted receipts $r$. Identity morphisms trivially satisfy the admissibility inequality $V(x) + 0 \le V(x) + 0$ and are accepted by $\mathrm{RV}$. For any composable pair of accepted morphisms $r_1: x_0 \to x_1$ and $r_2: x_1 \to x_2$, the syntactic composition $r_2 \circ r_1$ aggregates the spend and defect. By the Oplax Bound (Proposition 4.3), the composition is bounded. Addition of real values is strictly associative, ensuring $(r_3 \circ r_2) \circ r_1 = r_3 \circ (r_2 \circ r_1)$. Therefore, the structure forms a well-defined small category.
\end{proof}

% =====================================================
\section{Trace Algebra}

\begin{theorem}[Telescoping Trace Closure]
For any finite chain of admissible transitions $x_0 \xrightarrow{r_1} x_1 \xrightarrow{r_2} \dots \xrightarrow{r_n} x_n$, the global boundary satisfies:
\[
V(x_n)+\sum_{i=1}^n s_i \le V(x_0)+\sum_{i=1}^n d_i
\]
\end{theorem}
\begin{proof}
We proceed by induction on $n$. 
For $n=1$, the claim is exactly the definition of an admissible transition:
$V(x_1) + s_1 \le V(x_0) + d_1$.

Assume the claim holds for some $k \ge 1$:
\[
V(x_k) + \sum_{i=1}^k s_i \le V(x_0) + \sum_{i=1}^k d_i
\]
For the $(k+1)$-th step, admissibility guarantees:
\[
V(x_{k+1}) + s_{k+1} \le V(x_k) + d_{k+1}
\]
Adding $\sum_{i=1}^k s_i$ to both sides of this $(k+1)$-th inequality yields:
\[
V(x_{k+1}) + \sum_{i=1}^{k+1} s_i \le V(x_k) + d_{k+1} + \sum_{i=1}^k s_i
\]
Substituting $V(x_k) + \sum_{i=1}^k s_i$ using the inductive hypothesis:
\[
V(x_{k+1}) + \sum_{i=1}^{k+1} s_i \le \left( V(x_0) + \sum_{i=1}^k d_i \right) + d_{k+1}
\]
\[
V(x_{k+1}) + \sum_{i=1}^{k+1} s_i \le V(x_0) + \sum_{i=1}^{k+1} d_i
\]
By induction, the telescoping closure holds for all finite $n$. Intermediate states are algebraically annihilated, preventing internal budget inflation.
\end{proof}

% =====================================================
\section{Locally Posetal 2-Category Structure}

\begin{definition}[Slack 2-Cell]
A 2-cell $f \xrightarrow{\Delta} g$ between parallel morphisms exists if they share the same trace and spend, but $g$ provides more defect capacity:
\[
g.d = f.d + \Delta, \quad \Delta \ge 0
\]
\end{definition}

\begin{proposition}
$\mathbf{Coh}(\mathcal{S})$ forms a locally posetal 2-category.
\end{proposition}
\begin{proof}
The hom-sets $\mathbf{Coh}(x, y)$ are preordered by the defect bound. Since equality of morphisms is defined by the equality of their $(r, s, d)$ triplets, this preorder is a partial order. Vertical composition is addition of slack, and horizontal composition is additive whiskering. This makes the Coh category a 2-category where all 2-cells between parallel 1-cells are unique.
\end{proof}

\begin{theorem}[Semantic Initial Property]
For any observable trace $R$, the tightest possible receipt (Semantic Bound) is the initial object in the poset of valid receipts for that trace.
\end{theorem}
\begin{proof}
By Definition 4.2, $\delta(R)$ is the infimum of valid defect bounds for trace $R$. Thus, for any receipt $f$ with trace $R$, $\delta(R) \le f.d$. There exists a unique $\Delta = f.d - \delta(R) \ge 0$ defining a 2-cell from the semantic bound to $f$.
\end{proof}

% =====================================================
\section{Separation Law}

\begin{axiom}[Structural Independence]
The spend functional and the defect allowance are structurally independent:
\[
s(r) \not\equiv d(r)
\]
This guarantees that the thermodynamic penalty of execution cannot be trivially zeroed out by manipulating the defect reporting.
\end{axiom}

% =====================================================
\section{Worked Example: Discrete Trace}

Consider a simple cognitive state progression tracking a violation functional $V$:
\[
x_0 \xrightarrow{r_1} x_1 \xrightarrow{r_2} x_2
\]
Assume $V(x_0) = 10$. For $r_1$, the verifier yields $s_1=2, d_1=1$, so $V(x_1) \le 9$. For $r_2$, we have $s_2=3, d_2=1$, so $V(x_2) \le 7$. 
The global telescoping check:
\[
V(x_2) + (2 + 3) \le 10 + (1 + 1) \implies V(x_2) \le 7
\]
This demonstrates that evaluating the endpoints securely bounds the global behavior, regardless of internal state complexity.

% =====================================================
\section{Conclusion}

The Coh Category provides a rigorous foundation for governed computation by isolating physical constraints into a verifiable categorical ledger. By separating syntactic execution from semantic bounds via a 2-categorical structure, we achieve both operational strictness and physical soundness.

% =====================================================
\begin{thebibliography}{9}
\bibitem{kelly} Kelly, G. (1982). Basic Concepts of Enriched Category Theory.
\bibitem{benabou} Bénabou, J. (1967). Introduction to Bicategories.
\bibitem{coh} NoeticanLabs (2026). Coh Framework Specification v2.0.
\end{thebibliography}

\end{document}
